\documentclass[11pt,a4paper]{article}
\usepackage[margin=0.85in]{geometry}
\usepackage{booktabs}
\usepackage{amsmath}
\usepackage{hyperref}
\hypersetup{colorlinks=true,urlcolor=blue,linkcolor=black}
\title{A Human-Validated, Multilingual Research Design for CDP Climate Narratives}
\author{}
\date{29 September 2026}
\begin{document}
\maketitle

\section{Recommendation and scope}
The defensible contribution is a validated measurement framework for
\emph{reported} corporate climate actions, not a claim that an unsupervised
cluster identifies actual abatement or a causal effect. A suitable working
title is \emph{Beyond Questionnaire Categories: Human-Validated Multilingual
Measurement of Corporate Climate Actions in CDP Disclosures}.

Use 2024 Q7.55.2 initiative comments for initial taxonomy development and
validation. Extend the frozen taxonomy to initiative disclosures in
2016--2024 only after testing year comparability. Treat Q7.55.3 investment
mechanisms and Q7.55.4 explanations for no active initiatives as separate
2024 companion analyses. Risks and opportunities need separate codebooks:
a reported risk or opportunity is not an implemented initiative.

\paragraph{Research questions.}
(1) What action mechanisms do narratives describe beyond CDP's selected
categories? (2) Which unassigned responses express existing themes, multiple
themes, genuinely new themes, or insufficient information? (3) How sensitive
are assignments to language, translation, boilerplate, and text truncation?
(4) After validation, how does the prevalence of reported action categories
vary across firms and disclosure years?
Questions (1)--(3) are the primary measurement study; (4) is a descriptive
extension, not an automatic consequence of obtaining clusters.

\section{What the existing files establish}
The audit read the documentation, generating code, summary tables, selected
representative and outlier texts, all three 2024 text-assignment tables, and
selected columns across all rows of the currency-enriched longitudinal
hybrid table. It is not a human reading of every narrative, a complete
language census, or a refit of the models. Existing source files and model
assignments were left unchanged.

\begin{center}
\small
\begin{tabular}{lrrrr}
\toprule
2024 question & Rows & Companies & Topics & Outliers \\
\midrule
Q7.55.2 initiatives & 10,545 & 3,455 & 54 & 4,348 (41.23\%) \\
Q7.55.3 investment methods & 11,695 & 4,283 & 48 & 5,399 (46.17\%) \\
Q7.55.4 no-active explanations & 580 & 580 & 9 & 202 (34.83\%) \\
\bottomrule
\end{tabular}
\end{center}
Topic counts exclude $-1$; topic IDs are question-specific. Q7.55.1 has
4,967 complete company stage profiles grouped numerically, not by text.
Q7.56 has eight eligible descriptions and should be handled as cases rather
than a standalone clustering exercise.

The longitudinal enriched hybrid table contains 314,298 unique record IDs:
108,815 initiatives, 112,913 risks, and 92,570 opportunities. Only 18,000
rows have direct MiniLM embedding assignments; 294,907 have TF--IDF/SVM
surrogate assignments. Another 1,391 short narratives are unassigned.
The surrogate holdout macro-F1 values (0.752 initiatives, 0.624 risks,
0.671 opportunities) measure reproduction of model assignments, not
agreement with independently validated substantive labels.

\paragraph{Selection and dependence.}
The longitudinal extract has 18,397 initiative rows in 2024, from 5,413
companies. Only 10,545 comments meet the 80-character text-model threshold:
7,852 rows, or 42.68\%, are outside that fit. Of all initiative comments,
16,792 are nonblank. Audit excluded short and blank comments separately;
absence from the fitted text corpus is not absence of action.
Character thresholds can also select languages differently.

After case-folding and collapsing whitespace, there are 43,034 excess
duplicate narratives within response type in the longitudinal table.
Duplicates can be genuine repeated disclosures, so preserve source rows
and multiplicity while preventing train/test leakage. In 2024 Q7.55.2
topic 32, 37 of 48 rows come from one company (British Land); the remaining
rows come from two companies. This is a concentration warning, not by
itself proof of an invalid theme.

\paragraph{Language and content.}
Some topic words are language-driven, but their clusters are not reliable
language labels. For example, Q7.55.4 topic 0 contains both a Spanish
planning response and a Korean resource-constraint response. Outliers also
include clear renewable-electricity procurement and funding narratives.
Do not label all $-1$ records as ``no action'', ``other'', or unusable text.

\section{Construct definition and an English codebook}
The source response row is the observation unit. Allow several substantive
codes per row, with evidence spans identifying each action or barrier.
Do not convert several clauses into independent firm observations.
Keep five dimensions separate:
\begin{enumerate}
\item Action or mechanism: building/process efficiency, electrification,
renewable generation, renewable-electricity procurement and certificates,
transport and logistics, circularity, fugitive/process emissions,
supplier/value-chain action, and organizational/behavioral change.
\item Status: implemented, ongoing, planned, explicitly absent, or unclear.
Keep questionnaire-derived status separate from narrative-derived status.
\item Scope and target: use only explicit evidence or separately recorded
structured selections; do not infer Scope 3 from a vague supply-chain word.
\item Evidence specificity: named activity, implementation detail, quantified
reported estimate, general aspiration, or insufficient information.
These are disclosure attributes, not external assurance ratings.
\item Quality flags: multiple themes, translation uncertainty, boilerplate,
structured/narrative contradiction, or insufficient explanation.
\end{enumerate}
Use a distinct mechanism taxonomy for Q7.55.3 (budgets, financial appraisal,
carbon pricing, incentives, regulation, engagement, and partnerships).
Use a distinct barrier/explanation taxonomy for Q7.55.4 (resources,
measurement/planning, corporate dependencies, perceived limited potential,
timing, and insufficient explanation). Low materiality or technical
infeasibility must not be inferred merely from industry or company size.

Every code needs a definition, inclusion/exclusion rules, positive and
negative examples, and an evidence requirement. Record ``not stated''
separately from an explicit negative. Permit ``new theme'' and ``unclear''
instead of forcing a label. The pilot codes supplied with this note are
draft examples, not a frozen ontology.

\section{Sampling and human validation}
\paragraph{Development, not evaluation.}
Begin with approximately 60 calibration responses spanning assigned
topics, outliers, short answers, languages, and contradictions. Two human
coders independently read the originals or reviewed translations, discuss
disagreements, and revise the codebook. Then inspect random responses from
each cluster, not only centroid-nearest representatives. Representatives
are useful for discovery but tend to hide difficult cases.

The existing \texttt{prepare\_cdp\_annotation\_sample.py} provides a 330-row
lexical-cluster/year-balanced development sample across three response
types. It is useful prior work, but not a probability sample of 2024
BERTopic outliers or a blind evaluation set for this study. Its exposed
cluster labels must be hidden for independent coding.

\paragraph{Locked evaluation.}
Freeze the codebook, cluster-to-code mapping procedure, translation rules,
and tuning decisions before examining evaluation labels. An initial
budget is 400 Q7.55.2 and 200 Q7.55.3 evaluation rows, stratified by
assigned/outlier status and language, with explicit coverage of small
topics. These are planning budgets, not a power calculation or a guarantee
of precision for each subgroup. Around 385 independent simple-random
observations give a worst-case 95\% proportion margin near five percentage
points; stratification, company dependence, and small subgroups require a
design-specific calculation. Expand samples based on desired precision.

All 580 eligible Q7.55.4 responses are small enough for a human-coding
census; audit shorter excluded responses separately. Rows used to develop
the codebook or inspect examples remain development rows, even in a census,
and cannot also provide untouched test evidence.

Keep both random/probability evaluation and purposive diagnostic samples.
Record stratum sizes, sampling fractions, selection seed, and inclusion
probabilities. Use inverse-probability weights where appropriate for
population estimates; do not calculate prevalence directly from an
outlier-oversampled review queue. Audit short-answer exclusions against
the original source population rather than only fitted assignment files.

\paragraph{Blinding and independence.}
Human coders should first see the question, text, and necessary context,
not the machine topic, suggested code, or the other coder's answer. Where
feasible hide organization identity too. The researcher can be one coder;
AI proposals are not an independent human second coder. Record initial
codes before adjudication and preserve disagreement histories. Double-code
the locked evaluation sample; bilingual reviewers assess translation
fidelity for the language subsets they can actually read.

\paragraph{Metrics.}
For each substantive binary code, report coder agreement, confusion counts,
and an appropriate chance-corrected measure such as nominal Krippendorff's
alpha, with uncertainty and code prevalence. Do not apply a single nominal
agreement statistic blindly to unordered multi-label sets. After
adjudication, evaluate mapped model labels using per-code precision/recall
and macro/micro F1, together with abstention and coverage.
Include outliers in whole-system accounting; otherwise conditional accuracy
can look better simply by rejecting hard cases. Report metrics by language,
year, assignment method, text length, and questionnaire version, with
sample sizes and uncertainty. Use company-level resampling while respecting
the sampling strata.

\section{Translation and outlier handling}
Detect language on free-text segments rather than on concatenated English
CDP categories. Automatic language predictions are screening aids: retain
mixed/unknown labels and manually confirm uncertain cases.
Store the original, language, complete English translation, translator/model
and version, date, reviewer, uncertainty notes, and coding evidence.
Preserve negation, units, numerical claims, time references, and distinctions
between planned and implemented actions. A summary is not a translation.
Do not silently repair inconsistent source claims.

Compare (A) multilingual embeddings of original responses with
(B) embeddings of reviewed English translations on the same paired sample.
Compare both to bilingual substantive coding and inspect per-language
errors; English translation is not automatically ground truth.
Translation should not be assumed to solve language clustering without
this paired comparison. Use approved/local workflows for bulk processing;
no company response was sent to an external translation service in this audit.

For each outlier, record whether it fits an existing validated code, needs
several codes, suggests a new theme, contains insufficient information, or
requires language/data-quality review. These flags may coexist.
Only create a new category after reviewing repeated substantive evidence
across independent companies; preserve rare but meaningful cases without
requiring them to form a density cluster.

Keep \texttt{original\_cluster}, proposed assignments, final human codes,
assignment method, and review status as separate fields. Never overwrite
$-1$ with an unqualified success label. BERTopic's outlier reduction is a
reassignment mechanism, not validation \cite{bertopic}.
Tune reassignment similarity thresholds on development labels, then report
precision versus coverage on untouched responses. Retain abstention.
Scores from direct embedding cosine similarity and SVM margins are neither
interchangeable nor calibrated correctness probabilities.

\section{Model comparisons and longitudinal inference}
Use matched text populations for the following comparisons: narrative-only
versus structured-category-plus-narrative inputs; original multilingual
versus translated English text; lexical NMF versus embedding clustering;
and different seeds, density settings, and long-text handling.
CDP categories can inform the codebook, but agreement with those categories
is not independent evidence of discovery when they were input features.
Compare with a structured-category-only baseline to establish added value.

For long responses, measure tokenizer truncation using the actual saved
encoder configuration. The longitudinal hybrid metadata records 256 tokens;
the 2024 clustering script does not explicitly set a sequence limit.
Do not assume the same effective limit for both. Compare truncation with
sentence/chunk-based encoding on the same cases, preserving row identity.
Assess stability on a fixed anchor set across refits, aligning topics or
using permutation-invariant measures and reporting noise separately.
Stability and word coherence support, but cannot replace, human validity
\cite{grimmer}.

Split companies and exact/near-duplicate text groups before fitting
embeddings-to-cluster mappings, surrogates, or supervised classifiers used
for held-out evaluation. If duplicate text links several companies, keep
the connected group together or explicitly isolate it from evaluation.
For temporal generalization, learn the taxonomy/model on earlier years and
evaluate on later years without retuning. The existing full-period hybrid
fit cannot serve as an untouched prospective 2024 test.

Check the field-mapping audit and questionnaire glossary, especially the
2018 and 2024 changes. Do not equate disclosure year with implementation or
financial year. Compare repeated-company panels with the changing
cross-section and report selection/attrition. Use stable substantive code
IDs across years; independent annual topic numbers are not comparable.

For company-year code prevalence, use an indicator that at least one
eligible response expresses code $k$, and report the denominator of
companies with observable responses. Keep no disclosure, excluded/short
text, unresolved coding, and an explicit absence distinct. Report
response-level proportions separately. Do not sum multi-label percentages
as if categories were exclusive or rely only on a dominant-topic field:
that suppresses secondary themes and unresolved responses.

Any association with reported savings, costs, sector, or later emissions
requires additional design and measurement checks. Reported annual savings
are not verified emissions reductions; currencies must be harmonized
before monetary comparisons, and savings may overlap across projects.
Neither successful clustering nor company/year controls alone establish
causality.

\section{Translation-and-coding pilot supplied}
The complete-text field-separated implementation and human coding workflow
are documented in \path{reports/cdp_text_clustering/cdp_field_chunking_workflow.tex}.
It preserves exact evidence spans, pools chunks per response/source field,
and compares that representation with truncation of the same field.
Machine code suggestions remain distinct from accepted human evidence.

An additional empirical comparison is documented in
\path{reports/cdp_text_clustering/cdp_clustering_method_verification.tex}, with artifacts in
\path{data/outputs/cdp_text_clustering/cdp_clustering_verification_20260929/}.
It compares NMF, MiniLM/k-means, and the UMAP/HDBSCAN clustering core on
company/exact-text-grouped holdouts across five seeds. It supports a
conditional exploratory method choice, not independent human-label
validity or a state-of-the-art encoder claim.

The companion file
\path{data/outputs/cdp_text_clustering/cdp_2024_emissions_initiative_clusters/translation_coding_pilot.csv}
contains six non-English responses with complete draft English translations
and two English outlier examples. Each retains source filename, organization
ID, question, row order, original text, and original cluster.
It supplies proposed codes and empty human-review fields. All eight rows
are deliberately selected development examples, marked \texttt{unreviewed};
they are neither representative prevalence evidence nor a validated
reference set. The AI-generated code proposals should be hidden in any
independent human-coding exercise. Source assignments were not changed.

\begin{thebibliography}{9}
\bibitem{grimmer}
Grimmer, J. and Stewart, B. M. (2013).
Text as Data: The Promise and Pitfalls of Automatic Content Analysis Methods
for Political Texts. \emph{Political Analysis}, 21(3), 267--297.
\url{https://doi.org/10.1093/pan/mps028}.
\bibitem{bertopic}
BERTopic documentation. Outlier reduction.
\url{https://maartengr.github.io/BERTopic/getting_started/outlier_reduction/outlier_reduction.html}.
\end{thebibliography}
\end{document}
